\subsection{Adaptive Exploration Mechanism}
\label{sec:adaptive_epsilon_decay}

In standard time-driven MASA-QMIX style formulations, decision making is aligned
with a fixed decision-time grid (Key Property~1; Section~\ref{sec:standart_masaproblem_formulation}).
As a consequence, \emph{each episode contains a fixed number of decision steps},
and each reinforcement learning step corresponds to exactly one decision point.
Under this assumption, a time-/step-based $\epsilon$-decay schedule provides a
\emph{predictable fraction of exploratory decisions}.

\paragraph{Step-based $\epsilon$ decay under fixed-length episodes (baseline case).}
Let an episode contain exactly $K$ decision steps (e.g., $K=50$), and let training
span $E$ episodes (e.g., $E=100$). Then the total number of decision points is
deterministic:
\begin{equation}
D_{\text{tot}} = E \cdot K.
\label{eq:baseline_total_decisions}
\end{equation}
A standard linear step-based decay is
\begin{equation}
\epsilon(k)
=
\max\!\left(
\epsilon_{\min},
\;
\epsilon_0 - \frac{k}{K_{\mathrm{anneal}}}
(\epsilon_0 - \epsilon_{\min})
\right),
\qquad k=0,1,\dots
\label{eq:step_based_epsilon}
\end{equation}
where $k$ is the RL step index and $K_{\mathrm{anneal}}$ is chosen as a fixed
fraction of the total steps. For example, setting
\begin{equation}
K_{\mathrm{anneal}} = \rho \, D_{\text{tot}}, \qquad \rho=0.6,
\label{eq:anneal_fraction_baseline}
\end{equation}
means that exploration is gradually reduced until approximately $60\%$ of all
decision points have been processed.

Because the mapping "one RL step = one decision point" holds and $D_{\text{tot}}$
is fixed, this schedule is meaningful: after $\rho D_{\text{tot}}$ decisions,
the policy has indeed experienced $\rho D_{\text{tot}}$ interaction opportunities.

\paragraph{Why this breaks in the proposed event-driven setting.}
In the dynamic setting studied in this thesis, the environment is \emph{event-driven}
and decision points occur only when an execution event happens (Section~\ref{sec:discrete_event}),
i.e., after job arrivals and operation completions. Therefore, the number of
decision points per episode is not fixed.

Let $D_e$ denote the number of decision points in episode $e$.
Then the total number of decision points across $E$ episodes is random:
\begin{equation}
D_{\text{tot}} = \sum_{e=1}^{E} D_e,
\qquad
D_e \ \text{varies across episodes}.
\label{eq:dynamic_total_decisions}
\end{equation}
This variability is caused by (i) stochastic job arrivals, (ii) heterogeneous
operation lengths and sequences, and (iii) asynchronous execution dynamics
(Sections~\ref{sec:stochastic_hetero_job_arrivals},
\ref{sec:asynchron_decision_variable},
\ref{sec:discrete_event}).

Crucially, in this event-driven setting there is no guaranteed proportionality
between the episode index (or elapsed time) and the number of collected decisions:
\begin{equation}
\sum_{e=1}^{\lfloor \rho E \rfloor} D_e
\not\approx
\rho \sum_{e=1}^{E} D_e
\quad \text{in general.}
\label{eq:no_fraction_guarantee}
\end{equation}
Hence, an "anneal over the first $\rho E$ episodes" or "anneal over the first
$\rho$ fraction of elapsed time" does \emph{not} guarantee that $\rho$ fraction
of decision experience has been explored.

\subsubsection{Formal Analysis of Episode-Based Decay Failure}

\paragraph{Theorem 1 (Exploration Coverage Failure under Episode-Based Decay).}
\label{thm:episode_decay_failure}
Under episode-based $\epsilon$-decay with annealing parameter $\rho \in (0,1)$
over the first $\lfloor \rho E \rfloor$ episodes, the explored decision fraction
$f_{\text{explore}} = D_{\text{early}}/D_{\text{tot}}$ is bounded by:
\begin{equation}
\frac{\rho D_{\min}}{\rho D_{\min} + (1-\rho) D_{\max}}
\leq f_{\text{explore}} \leq
\frac{\rho D_{\max}}{\rho D_{\max} + (1-\rho) D_{\min}},
\label{eq:exploration_bounds}
\end{equation}
where $D_{\min} = \min_{e \in [1,E]} D_e$ and $D_{\max} = \max_{e \in [1,E]} D_e$.
In general, $f_{\text{explore}} \neq \rho$, and the deviation can be arbitrarily
large when $D_{\max}/D_{\min}$ is large.

\paragraph{Proof.}
The worst-case lower bound occurs when the first $\lfloor \rho E \rfloor$ episodes
all have minimum decision count $D_{\min}$, and the remaining episodes have maximum
count $D_{\max}$:
\begin{equation}
f_{\text{explore}}^{\text{worst}}
=
\frac{\rho E D_{\min}}{\rho E D_{\min} + (1-\rho)E D_{\max}}
=
\frac{\rho D_{\min}}{\rho D_{\min} + (1-\rho) D_{\max}}.
\end{equation}
The upper bound occurs in the opposite case (early episodes with $D_{\max}$,
later with $D_{\min}$). The inequality $f_{\text{explore}} \neq \rho$ follows
whenever $D_{\min} \neq D_{\max}$. $\square$

\paragraph{Concrete illustrations.}
To demonstrate the severity of this problem, consider $E=100$ episodes with 
target exploration fraction $\rho=0.6$ (i.e., decay completes after episode 60):

\textbf{Extreme Scenario A (low decision density early, high later):}
If episodes 1--60 each have $D_e=10$ decision points while episodes 61--100 have
$D_e=40$, then Theorem~\ref{thm:episode_decay_failure} yields:
\begin{equation}
f_{\text{explore}}^{A}
=
\frac{0.6 \cdot 10}{0.6 \cdot 10 + 0.4 \cdot 40}
=
\frac{6}{22}
\approx 0.27.
\label{eq:scenario_A_fraction}
\end{equation}
Despite intending $60\%$ exploration, the policy explores only $27\%$ of decision
opportunities, severely underexploring the state space.

\textbf{Extreme Scenario B (high decision density early, low later):}
If episodes 1--60 each have $D_e=40$ while episodes 61--100 have $D_e=10$:
\begin{equation}
f_{\text{explore}}^{B}
=
\frac{0.6 \cdot 40}{0.6 \cdot 40 + 0.4 \cdot 10}
=
\frac{24}{28}
\approx 0.86.
\label{eq:scenario_B_fraction}
\end{equation}
Here, exploration persists through $86\%$ of decisions, wasting exploitation
potential after the policy has likely converged.

Therefore, the \emph{same} episode-based decay rule produces coverage ranging
from $27\%$ to $86\%$ -- a $59$ percentage point difference -- depending purely
on when decision-intensive episodes occur.

This demonstrates the core problem: in an event-driven environment, annealing as
a function of time or episode index cannot guarantee a predictable amount of
exploratory decision experience. The plan "explore for the first $60\%$ of
training" is not well-defined in terms of actual interaction data.

\subsubsection{Decision-Point-Based Exploration Decay with Guaranteed Coverage}

\paragraph{Proposed formulation.}
To ensure that exploration decreases only after a controlled amount of decision
experience has been collected, the proposed framework defines $\epsilon$ as a
function of the cumulative decision point index $d$:
\begin{equation}
\epsilon_d
=
\max\!\left(
\epsilon_{\min},
\;
\epsilon_0 - \frac{d}{D_{\mathrm{anneal}}}
(\epsilon_0 - \epsilon_{\min})
\right),
\qquad d=0,1,\dots
\label{eq:decision_based_epsilon_final}
\end{equation}
where $d$ counts actual decision points encountered across episodes.
By choosing $D_{\mathrm{anneal}}$ directly (e.g., "anneal over the first
$D_{\mathrm{anneal}}=25{,}000$ decision points"), the schedule guarantees that the
policy explores over a controlled number of interaction opportunities
independent of how decision points are distributed across episodes.

\paragraph{Theorem 2 (Guaranteed Exploration Coverage under Decision-Based Decay).}
\label{thm:decision_based_guarantee}
Under decision-point-based $\epsilon$-decay with annealing horizon
$D_{\mathrm{anneal}}$, the number of exploratory decisions satisfies:
\begin{equation}
N_{\text{explore}}
= \int_0^{D_{\mathrm{anneal}}} \epsilon_d \, \mathrm{d}d
= \frac{D_{\mathrm{anneal}}(\epsilon_0 + \epsilon_{\min})}{2},
\label{eq:guaranteed_exploration}
\end{equation}
independent of the episode-level distribution of decision points. Furthermore,
the explored decision fraction is:
\begin{equation}
f_{\text{explore}}
= \frac{N_{\text{explore}}}{D_{\text{tot}}}
= \frac{D_{\mathrm{anneal}}}{D_{\text{tot}}} \cdot
\frac{\epsilon_0 + \epsilon_{\min}}{2},
\label{eq:guaranteed_fraction}
\end{equation}
which is deterministic given $D_{\mathrm{anneal}}$ and $D_{\text{tot}}$.

\paragraph{Proof.}
Under linear decay (Eq.~\ref{eq:decision_based_epsilon_final}), the exploration
rate decreases linearly from $\epsilon_0$ to $\epsilon_{\min}$ over
$D_{\mathrm{anneal}}$ decision points. Assuming $\epsilon$-greedy action selection,
the expected number of exploratory actions per decision is $\epsilon_d$. Integrating:
\begin{align}
N_{\text{explore}}
&= \int_0^{D_{\mathrm{anneal}}} \epsilon_d \, \mathrm{d}d
= \int_0^{D_{\mathrm{anneal}}}
\left(
\epsilon_0 - \frac{d}{D_{\mathrm{anneal}}}(\epsilon_0 - \epsilon_{\min})
\right) \mathrm{d}d
\\
&= \epsilon_0 D_{\mathrm{anneal}}
- \frac{\epsilon_0 - \epsilon_{\min}}{D_{\mathrm{anneal}}}
\cdot \frac{D_{\mathrm{anneal}}^2}{2}
= D_{\mathrm{anneal}}
\left(
\epsilon_0 - \frac{\epsilon_0 - \epsilon_{\min}}{2}
\right)
\\
&= \frac{D_{\mathrm{anneal}}(\epsilon_0 + \epsilon_{\min})}{2}.
\end{align}
This quantity depends only on $D_{\mathrm{anneal}}$, $\epsilon_0$, and
$\epsilon_{\min}$, not on how decision points are distributed across episodes. $\square$

\paragraph{Proposition 1 (Exploration Variance Elimination).}
\label{prop:variance_elimination}
Let $D_e \sim \mathcal{N}(\mu_D, \sigma_D^2)$ be i.i.d.\ episode decision counts.
Under episode-based decay with target fraction $\rho$, the variance of the number
of exploratory decisions is:
\begin{equation}
\mathrm{Var}\!\left[N_{\text{explore}}^{\text{episode}}\right]
= \bar{\epsilon}^2 \cdot \lfloor \rho E \rfloor \cdot \sigma_D^2,
\label{eq:variance_episode}
\end{equation}
where $\bar{\epsilon} = (\epsilon_0 + \epsilon_{\min})/2$ is the average
exploration rate. Under decision-based decay:
\begin{equation}
\mathrm{Var}\!\left[N_{\text{explore}}^{\text{decision}}\right] = 0.
\label{eq:variance_decision}
\end{equation}

\paragraph{Proof.}
Under episode-based decay, the number of exploratory decisions is approximately
$N_{\text{explore}}^{\text{episode}} \approx \bar{\epsilon} \sum_{e=1}^{\lfloor \rho E \rfloor} D_e$.
Since $D_e$ are i.i.d.\ with variance $\sigma_D^2$:
\begin{equation}
\mathrm{Var}\!\left[\sum_{e=1}^{\lfloor \rho E \rfloor} D_e\right]
= \lfloor \rho E \rfloor \sigma_D^2,
\end{equation}
and thus
$\mathrm{Var}[N_{\text{explore}}^{\text{episode}}]
= \bar{\epsilon}^2 \lfloor \rho E \rfloor \sigma_D^2$.
Under decision-based decay, $N_{\text{explore}}^{\text{decision}}$ is deterministic
(Theorem~\ref{thm:decision_based_guarantee}), so its variance is zero. $\square$

\paragraph{Practical implications.}
In the experimental setting (Section~\ref{sec:experiments}), observed episode
decision counts range from $D_e \in [37, 51]$ with $\bar{D}=41.66$ and
$\sigma_D=2.70$, confirming such variability is not hypothetical but occurs in
practice. Setting $D_{\mathrm{anneal}}=25{,}000$ with $\epsilon_0=1.0$ and
$\epsilon_{\min}=0.05$ guarantees:
\begin{equation}
N_{\text{explore}}
= \frac{25{,}000 \times (1.0 + 0.05)}{2}
= 13{,}125
\text{ exploratory decisions},
\label{eq:practical_example}
\end{equation}
regardless of how these decision points are distributed across episodes.

Consequently, the proposed adaptive exploration mechanism:
\begin{itemize}
\item \textbf{Aligns exploration with interaction structure:} Decay proceeds at
the rate agents actually interact with the environment.

\item \textbf{Prevents premature convergence:} Low-density episodes do not cause
premature exploitation.

\item \textbf{Ensures robust learning:} Consistent exploration experience across
runs with different stochastic realizations.

\item \textbf{Enables principled hyperparameter tuning:} $D_{\mathrm{anneal}}$ has
a clear semantic meaning (number of exploratory decisions), unlike
$E_{\mathrm{anneal}}$ (number of episodes, whose decision content varies).
\end{itemize}

This mechanism is essential for achieving stable and reproducible learning in
stochastic, event-driven multi-agent scheduling environments.
